在已知原评估结果后，我们对冻结的共同候选分配研究进行事后敏感性分析。在原最佳验证BER并列集合内（容差$10^{-12}$），分别选择最小与最大关系系数，改变了300次选择中的86次与100次。原默认优先结果仍为主要结果。在20项独立选参ACC比较中，两种替代规则分别有18项与16项Full平均差为正，原规则为19项；采用各政策Full所选系数的控制分别有19项与19项为正，原规则为19项。三套政策保留相同数据、查询流、有效折和系数网格。这些结果描述既有语料内对并列取舍的依赖，不构成独立确认，也不用于根据测试分数选择有利政策。
19/20项正均值这一精确计数依赖并列规则。不过，三套规则下独立选参的ACC均值对比均以正差为多数。这些符号计数不确立逐运行、多指标或未经测试规则下的普遍改善。
